% ============================================================
%  AFD — Validação de Links da Wikipédia (pt.wikipedia.org)
%  Linguagens Formais e Compiladores — 6º Período — 2026
%
%  Como compilar:
%    pdflatex afd_wikipedia.tex
%  Ou abra no Overleaf (overleaf.com) e compile lá.
% ============================================================

\documentclass[a4paper, 12pt]{article}

\usepackage[utf8]{inputenc}
\usepackage[T1]{fontenc}
\usepackage[portuguese]{babel}
\usepackage{amsmath, amssymb}
\usepackage{xcolor}
\usepackage{geometry}
\usepackage{tikz}
\usetikzlibrary{automata, positioning, arrows.meta, fit, backgrounds}
\usepackage{booktabs}
\usepackage{caption}
\usepackage{pdflscape}   % para páginas em landscape dentro do doc

% ── Geometria padrão (retrato) ───────────────────────────────
\geometry{a4paper, margin=2cm}

% ── Paleta de cores ──────────────────────────────────────────
\definecolor{Cblue}{HTML}{1A73E8}   % estados normais
\definecolor{Cgreen}{HTML}{188038}  % estado de aceitação
\definecolor{Cred}{HTML}{D93025}    % estado poço (trap)
\definecolor{Cgray}{HTML}{5F6368}   % setas / texto secundário

% ── Estilos TikZ ─────────────────────────────────────────────
\tikzset{
  % Estado normal
  every state/.style={
    circle, draw=Cblue, fill=Cblue!8, thick,
    minimum size=1cm, inner sep=0pt,
    font=\small\bfseries, text=Cblue
  },
  % Estado de aceitação (sobrescreve a cor)
  accepting/.style={
    draw=Cgreen, fill=Cgreen!10, text=Cgreen,
    double, double distance=2.5pt
  },
  % Estado poço / trap
  trap/.style={
    draw=Cred, fill=Cred!8, text=Cred, dashed
  },
  % Setas padrão
  every edge/.style={
    draw=Cgray, ->, >=stealth, shorten >=1pt, shorten <=1pt,
    font=\scriptsize
  }
}

% ── Título ───────────────────────────────────────────────────
\title{\textbf{Autômato Finito Determinístico}\\[4pt]
       \large Validação de Links da Wikipédia em Português\\[2pt]
       \normalsize Projeto 1 — Linguagens Formais e Compiladores}
\author{Ciência da Computação · 6º Período}
\date{2026}

% ============================================================
\begin{document}
\maketitle
\thispagestyle{empty}

% ── Descrição do problema ────────────────────────────────────
\section*{Linguagem Reconhecida}

O AFD reconhece exatamente a linguagem descrita pela expressão regular:

\[
  L = \big\{\, w \in \Sigma^* \;\big|\;
  w \text{ casa com }
  \underbrace{\texttt{\^{}https://pt\backslash.wikipedia\backslash.org/wiki/}}_{\text{prefixo fixo — 30 símbolos}}
  \underbrace{\texttt{[A-Za-z0-9\_(){\%}-]+}}_{\text{nome do artigo ($\geq 1$ char)}}
  \texttt{\$}
  \,\big\}
\]

\subsection*{Definição Formal}

\[
  M = (Q,\ \Sigma,\ \delta,\ q_0,\ F)
\]

\[
  Q = \{q_0,\; q_1,\; \ldots,\; q_{30},\; q_{31},\; q_{\text{erro}}\}
  \quad (33 \text{ estados})
\]

\[
  \Sigma = \text{caracteres ASCII imprimíveis}
\]

\[
  q_0 = q_0 \qquad (\text{estado inicial})
\]

\[
  F = \{q_{31}\} \qquad (\text{único estado de aceitação})
\]

\bigskip
\noindent\textbf{Função de transição} $\delta$ :

\[
  \delta(q_i,\;\texttt{PREFIXO}[i]) = q_{i+1},
  \qquad 0 \leq i \leq 29
\]
\[
  \delta(q_i,\; c \neq \texttt{PREFIXO}[i]) = q_{\text{erro}},
  \qquad 0 \leq i \leq 29
\]
\[
  \delta(q_{30},\; c) =
  \begin{cases}
    q_{31}        & \text{se } c \in [A\text{-}Za\text{-}z0\text{-}9\_(){\%}\text{-}] \\
    q_{\text{erro}} & \text{caso contrário}
  \end{cases}
\]
\[
  \delta(q_{31},\; c) =
  \begin{cases}
    q_{31}        & \text{se } c \in [A\text{-}Za\text{-}z0\text{-}9\_(){\%}\text{-}] \\
    q_{\text{erro}} & \text{caso contrário}
  \end{cases}
\]
\[
  \delta(q_{\text{erro}},\; c) = q_{\text{erro}},
  \qquad \forall\, c \in \Sigma \quad (\text{estado poço})
\]

\smallskip
\noindent onde \; $\texttt{PREFIXO} = \texttt{"https://pt.wikipedia.org/wiki/"}$
\; (30 símbolos).

% ============================================================
%  DIAGRAMA 1 — Visão Geral (página em retrato)
% ============================================================
\newpage
\section*{Diagrama 1 — Visão Geral (Abstrata)}

\begin{figure}[h]
\centering
\begin{tikzpicture}[node distance=3.2cm and 3.2cm]

  % ── Nós ───────────────────────────────────────────────────
  \node[state, initial above]  (q0)  {$q_0$};

  % Grupo de estados do prefixo (retângulo arredondado)
  \node[draw=Cblue, fill=Cblue!8, thick, rounded corners=8pt,
        minimum width=2.8cm, minimum height=1.1cm,
        right=3.5cm of q0,
        font=\small\bfseries, text=Cblue]
        (pref) {$q_1 \;\cdots\; q_{29}$};

  \node[state, right=3.5cm of pref]           (q30) {$q_{30}$};
  \node[state, accepting, right=3.5cm of q30] (q31) {$q_{31}$};
  \node[state, trap, below=2.5cm of q30]      (err) {$q_{\text{erro}}$};

  % ── Transições principais ──────────────────────────────────
  \draw (q0)   edge[above] node{\texttt{h,t,t,p,s,:,/,/,p,t,...}} (pref);
  \draw (pref) edge[above] node{$\texttt{/}$ \quad (30\textdegree{} símbolo — prefixo completo)} (q30);
  \draw (q30)  edge[above] node{char válido} (q31);
  \draw (q31)  edge[loop above, looseness=6] node{char válido} (q31);

  % ── Transições de erro ────────────────────────────────────
  \draw (q0)   edge[bend right=35, below left]  node{inválido} (err);
  \draw (pref) edge[below, sloped]              node{símbolo errado} (err);
  \draw (q30)  edge[right]                      node{inválido} (err);
  \draw (q31)  edge[bend left=35, below right]  node{inválido} (err);
  \draw (err)  edge[loop below, looseness=6]    node{qualquer} (err);

  % ── Rótulos explicativos ──────────────────────────────────
  \node[below=0.15cm of q31, font=\tiny\bfseries, text=Cgreen] {\textsc{aceita}};
  \node[below=0.15cm of err, font=\tiny\bfseries, text=Cred]   {\textsc{trap}};

\end{tikzpicture}
\caption{Visão abstrata do AFD. Os estados $q_1 \cdots q_{29}$ representam
         a leitura símbolo a símbolo do prefixo fixo
         \texttt{https://pt.wikipedia.org/wiki/}.
         \textbf{char válido} $= [A\text{-}Za\text{-}z0\text{-}9\_(){\%}\text{-}]$.}
\end{figure}

% ============================================================
%  DIAGRAMA 2 — AFD Completo em Serpentina (página landscape)
% ============================================================
\begin{landscape}
\section*{Diagrama 2 — AFD Completo (todos os 33 estados)}

\begin{figure}[h]
\centering
\begin{tikzpicture}[node distance=1.55cm and 1.55cm]

  % ── Fileira 1: q0 → q10  (lê "https://pt") ───────────────
  \node[state, initial above]     (q0)  {$q_0$};
  \node[state, right=of q0]       (q1)  {$q_1$};
  \node[state, right=of q1]       (q2)  {$q_2$};
  \node[state, right=of q2]       (q3)  {$q_3$};
  \node[state, right=of q3]       (q4)  {$q_4$};
  \node[state, right=of q4]       (q5)  {$q_5$};
  \node[state, right=of q5]       (q6)  {$q_6$};
  \node[state, right=of q6]       (q7)  {$q_7$};
  \node[state, right=of q7]       (q8)  {$q_8$};
  \node[state, right=of q8]       (q9)  {$q_9$};
  \node[state, right=of q9]       (q10) {$q_{10}$};

  % ── Fileira 2: q11 → q20  (lê ".wikipedia") ──────────────
  \node[state, below=2cm of q10]  (q11) {$q_{11}$};
  \node[state, left=of q11]       (q12) {$q_{12}$};
  \node[state, left=of q12]       (q13) {$q_{13}$};
  \node[state, left=of q13]       (q14) {$q_{14}$};
  \node[state, left=of q14]       (q15) {$q_{15}$};
  \node[state, left=of q15]       (q16) {$q_{16}$};
  \node[state, left=of q16]       (q17) {$q_{17}$};
  \node[state, left=of q17]       (q18) {$q_{18}$};
  \node[state, left=of q18]       (q19) {$q_{19}$};
  \node[state, left=of q19]       (q20) {$q_{20}$};

  % ── Fileira 3: q21 → q30  (lê ".org/wiki/") ──────────────
  \node[state, below=2cm of q20]  (q21) {$q_{21}$};
  \node[state, right=of q21]      (q22) {$q_{22}$};
  \node[state, right=of q22]      (q23) {$q_{23}$};
  \node[state, right=of q23]      (q24) {$q_{24}$};
  \node[state, right=of q24]      (q25) {$q_{25}$};
  \node[state, right=of q25]      (q26) {$q_{26}$};
  \node[state, right=of q26]      (q27) {$q_{27}$};
  \node[state, right=of q27]      (q28) {$q_{28}$};
  \node[state, right=of q28]      (q29) {$q_{29}$};
  \node[state, right=of q29]      (q30) {$q_{30}$};

  % ── Fileira 4: q31 (aceitação) e q_erro (trap) ───────────
  \node[state, accepting, below=2cm of q30] (q31) {$q_{31}$};
  \node[state, trap, left=5cm of q31]       (err) {$q_{\text{e}}$};

  % ── Setas: Fileira 1 ─────────────────────────────────────
  \draw (q0)  edge[above] node{\texttt{h}} (q1);
  \draw (q1)  edge[above] node{\texttt{t}} (q2);
  \draw (q2)  edge[above] node{\texttt{t}} (q3);
  \draw (q3)  edge[above] node{\texttt{p}} (q4);
  \draw (q4)  edge[above] node{\texttt{s}} (q5);
  \draw (q5)  edge[above] node{\texttt{:}} (q6);
  \draw (q6)  edge[above] node{\texttt{/}} (q7);
  \draw (q7)  edge[above] node{\texttt{/}} (q8);
  \draw (q8)  edge[above] node{\texttt{p}} (q9);
  \draw (q9)  edge[above] node{\texttt{t}} (q10);

  % ── Curva: Fileira 1 → Fileira 2 ─────────────────────────
  \draw (q10) edge[right] node{\texttt{.}} (q11);

  % ── Setas: Fileira 2 ─────────────────────────────────────
  \draw (q11) edge[below] node{\texttt{w}} (q12);
  \draw (q12) edge[below] node{\texttt{i}} (q13);
  \draw (q13) edge[below] node{\texttt{k}} (q14);
  \draw (q14) edge[below] node{\texttt{i}} (q15);
  \draw (q15) edge[below] node{\texttt{p}} (q16);
  \draw (q16) edge[below] node{\texttt{e}} (q17);
  \draw (q17) edge[below] node{\texttt{d}} (q18);
  \draw (q18) edge[below] node{\texttt{i}} (q19);
  \draw (q19) edge[below] node{\texttt{a}} (q20);

  % ── Curva: Fileira 2 → Fileira 3 ─────────────────────────
  \draw (q20) edge[left] node{\texttt{.}} (q21);

  % ── Setas: Fileira 3 ─────────────────────────────────────
  \draw (q21) edge[above] node{\texttt{o}} (q22);
  \draw (q22) edge[above] node{\texttt{r}} (q23);
  \draw (q23) edge[above] node{\texttt{g}} (q24);
  \draw (q24) edge[above] node{\texttt{/}} (q25);
  \draw (q25) edge[above] node{\texttt{w}} (q26);
  \draw (q26) edge[above] node{\texttt{i}} (q27);
  \draw (q27) edge[above] node{\texttt{k}} (q28);
  \draw (q28) edge[above] node{\texttt{i}} (q29);
  \draw (q29) edge[above] node{\texttt{/}} (q30);

  % ── Curva: Fileira 3 → Fileira 4 ─────────────────────────
  \draw (q30) edge[right] node[align=left]{char\\válido} (q31);

  % ── Auto-laço q31 ────────────────────────────────────────
  \draw (q31) edge[loop right, looseness=5]
              node[right]{char válido} (q31);

  % ── Transições de erro (representativas) ─────────────────
  \draw (q0)  edge[bend right=25, below left] node{inválido} (err);
  \draw (q30) edge[above, sloped]             node{inválido} (err);
  \draw (q31) edge[above]                     node{inválido} (err);
  \draw (err) edge[loop left, looseness=6]    node{qualquer} (err);

  % ── Rótulos de grupo (anotações de fundo) ─────────────────
  \begin{scope}[on background layer]
    \node[draw=Cblue!30, fill=Cblue!5, rounded corners, inner sep=6pt,
          fit=(q0)(q10), label=above:{\color{Cgray}\scriptsize\textit{Lendo \texttt{"https://pt"}}}]
         {};
    \node[draw=Cblue!30, fill=Cblue!5, rounded corners, inner sep=6pt,
          fit=(q11)(q20), label=below:{\color{Cgray}\scriptsize\textit{Lendo \texttt{".wikipedia"}}}]
         {};
    \node[draw=Cblue!30, fill=Cblue!5, rounded corners, inner sep=6pt,
          fit=(q21)(q30), label=above:{\color{Cgray}\scriptsize\textit{Lendo \texttt{".org/wiki/"}}}]
         {};
  \end{scope}

  % ── Legenda ───────────────────────────────────────────────
  \node[below=0.2cm of q31, font=\tiny\bfseries, text=Cgreen] {\textsc{aceita}};
  \node[below=0.2cm of err, font=\tiny\bfseries, text=Cred]   {\textsc{trap}};

\end{tikzpicture}
\caption{AFD completo com todos os 33 estados em layout serpentina.
         Os estados estão agrupados em três fileiras correspondendo às
         três partes do prefixo fixo.
         \textbf{char válido} $= [A\text{-}Za\text{-}z0\text{-}9\_(){\%}\text{-}]$.
         Todas as transições indefinidas (não representadas) levam a
         $q_{\text{e}} = q_{\text{erro}}$ (estado poço).}
\end{figure}

\end{landscape}

% ============================================================
%  DIAGRAMA 3 — Zoom nos estados finais
% ============================================================
\newpage
\section*{Diagrama 3 — Estados Finais e Estado Poço}

\begin{figure}[h]
\centering
\begin{tikzpicture}[node distance=4cm]

  \node[state, draw=Cgray, fill=gray!8, text=Cgray]          (q30) {$q_{30}$};
  \node[state, accepting, right=of q30]                       (q31) {$q_{31}$};
  \node[state, trap, below=2.5cm of q30]                     (err) {$q_{\text{erro}}$};

  % Entrada em q30 (vinda do restante do AFD)
  \draw[dashed, Cgray] (-2.5,0) -- (q30)
        node[midway, above, font=\scriptsize, text=Cgray]
        {de $q_{29}$ após ler \texttt{'/'}};

  % q30 → q31 (1º char do artigo)
  \draw (q30) edge[above, bend left=15]
              node{$c \in [A\text{-}Za\text{-}z0\text{-}9\_(){\%}\text{-}]$} (q31);

  % q31 → q31 (auto-laço)
  \draw (q31) edge[loop right, looseness=5]
              node[right, align=left]{$c \in [A\text{-}Za\text{-}z0\text{-}9\_(){\%}\text{-}]$\\(auto-laço)}
              (q31);

  % q30 → erro
  \draw (q30) edge[right, bend right=10]
              node[right]{$c \notin [A\text{-}Za\text{-}z0\text{-}9\_(){\%}\text{-}]$} (err);

  % q31 → erro
  \draw (q31) edge[bend left=15, below right]
              node{$c \notin [A\text{-}Za\text{-}z0\text{-}9\_(){\%}\text{-}]$} (err);

  % err → err (auto-laço)
  \draw (err) edge[loop below, looseness=6]
              node{$\forall\, c \in \Sigma$} (err);

  % Rótulos
  \node[above=0.15cm of q30, font=\footnotesize, text=Cgray]
        {prefixo completo};
  \node[above=0.15cm of q31, font=\footnotesize\bfseries, text=Cgreen]
        {\textsc{estado de aceitação}};
  \node[below=0.15cm of err, font=\footnotesize\bfseries, text=Cred]
        {\textsc{estado poço (trap)}};

\end{tikzpicture}
\caption{Detalhe dos estados $q_{30}$, $q_{31}$ (aceitação) e $q_{\text{erro}}$ (poço).
         $q_{30}$ exige ao menos um caractere válido para passar a $q_{31}$,
         implementando o quantificador \texttt{+} da expressão regular.
         Uma vez em $q_{\text{erro}}$, nenhum símbolo é capaz de sair dele.}
\end{figure}

% ============================================================
%  Tabela de equivalência ER ↔ AFD
% ============================================================
\newpage
\section*{Equivalência ER $\;\longleftrightarrow\;$ AFD (Teorema de Kleene)}

\begin{table}[h]
\centering
\renewcommand{\arraystretch}{1.4}
\begin{tabular}{lll}
\toprule
\textbf{Fragmento da ER} & \textbf{Descrição} & \textbf{AFD equivalente} \\
\midrule
\texttt{\^{}} & Âncora de início &
  Estado inicial $q_0$ \\
\texttt{https://pt\textbackslash.wikipedia\textbackslash.org/wiki/} & Prefixo fixo (30 chars) &
  $\delta(q_i, \texttt{PREFIXO}[i]) = q_{i+1}$, $i=0..29$ \\
\texttt{[A-Za-z0-9\_(){\%}-]} & Classe de caracteres &
  Conjunto $C$ na função $\delta$ \\
\texttt{+} & Um ou mais & $q_{30} \xrightarrow{C} q_{31}$ (1\textordmasculine{} char) \\
           &             & $q_{31} \xrightarrow{C} q_{31}$ (demais) \\
\texttt{\$} & Âncora de fim & Aceitação apenas em $q_{31}$ após consumir toda a cadeia \\
qualquer outro char & Rejeição & $q_{\text{erro}}$ (estado poço) \\
\bottomrule
\end{tabular}
\caption{Mapeamento entre cada elemento da ER e os estados/transições do AFD.}
\end{table}

\bigskip

\noindent\textbf{Verificação no código Python:}

\medskip
\noindent Os dois métodos (AFD construído à mão e \texttt{re.match()} do Python)
são executados para o mesmo link e o resultado é comparado.
Por ser a mesma linguagem regular representada de formas diferentes,
eles \textbf{sempre concordam}, validando o Teorema de Kleene na prática.

\end{document}
